\section{Case Study: East African Montane Species}

\label{sec:case_study}
% ══════════════════════════════════════════════════════════════════════════════

We present a detailed case study of 43 montane bird species in the
Eastern Arc Mountains and Albertine Rift of East Africa. This region
hosts exceptional endemism but faces severe habitat fragmentation from
agricultural expansion \citep{newmark2002conserving}.

\subsection{Current Distribution Modeling}

HabitatNet identifies 12 distinct montane forest habitat patches, ranging
from 23\,km$^2$ (Nguru Mountains) to 4{,}800\,km$^2$ (Albertine Rift
complex). Species richness peaks at 1{,}800--2{,}400\,m elevation, with
sharp declines above 3{,}200\,m and below 1{,}200\,m.

\subsection{Climate-Driven Range Shifts}

Under SSP2-4.5 by 2081--2100, the model projects:
\begin{itemize}
  \item Mean upslope shift of the suitable climate envelope by
    340\,m ($\pm$120\,m).
  \item Loss of 38\% of currently suitable montane forest area due to
    ``summit trap'' effects---species near mountain peaks cannot
    shift further upward \citep{freeman2018climate}.
  \item Complete loss of suitable habitat for 4 narrow-endemic species
    with elevation ranges $<$500\,m.
\end{itemize}

\subsection{Priority Corridors}

The corridor algorithm identifies three critical connectivity zones:
\begin{enumerate}
  \item \textbf{Uluguru--Udzungwa link}: A 47\,km lowland gap currently
    acting as a dispersal barrier, where targeted restoration of 12\,km$^2$
    of riparian forest would connect 230\,km$^2$ of montane habitat.
  \item \textbf{Albertine Rift mid-elevation belt}: A 2{,}100\,m
    elevation corridor linking three national parks.
  \item \textbf{Usambara--Pare bridge}: Restoration of 8\,km$^2$ of
    degraded forest to reconnect two endemic bird areas.
\end{enumerate}


% ══════════════════════════════════════════════════════════════════════════════
